Zu Beginn der Projektarbeit haben wir, Wuyi und Arthur, uns zusammengesetzt und die Aufgaben, die uns als Gruppe \emph{Item- und Charakterdesign} zugeteilt wurden, klar definiert. Die Aufgaben sind folgende:

\begin{itemize}
	\item Festlegen der Anzahl an Lehrern und Items, sowie Auswählen der im Spiel vertretenen Lehrer für die jeweiligen Fachschaften
	\item Zeichnen der Spielerfigur sowie der Lehrer und im Spiel benutzten Items
	\item Implementieren des Spielers, der Lehrer und Items in das Programm 
\end{itemize}

\subsection{Spielerfigur}
Die Spielerfigur soll einerseits visuell den Benutzer im Spiel involvieren, dafür gibt es eine Methode, welche es ermöglicht, Haare, Hautfarbe und Kleidung individuell einzufärben.\\ Anderseits stellt sie den Hauptcharakter des Spieles dar und besitzt dementsprechend Funktionen, welche den Kampf gegen die Lehrer ermöglichen, wie auch ein Inventar, in welchem die erhaltenen Items gespeichert und mit Bild und Namen angezeigt werden. Dieses kann durch Anklicken der Spielerfigur geöffnet und mit einem Button wieder geschlossen werden kann. Die für den Kampf relevanten Funktionen wurden mit Absprache von der Gruppe \emph{Kampfkonzept} eingefügt und lagen außerhalb unserer Zuständigkeit.\\

\subsection{Lehrer*innen}
Die Lehrer und Lehrerinnen stellen in dem Spiel die Antagonisten dar, die es zu besiegen gilt. Geplant waren 3 Lehrer*innen für jede Fachschaft, mit variierenden Schwierigkeitsgraden im Kampf. Um diese darzustellen braucht es sowohl ein Bild des Lehrers/ der Lehrerin, welches dann gezeichnet wird, als auch Eigenschaften wie ein Name, der Schwierigkeitsgrad oder das "Level", die Zugehörigkeit zu einer Fachschaft, eine Menge an Lebenspunkten oder "Healthpoints" und eine Auswahl von Items, welche im Kampf zur Verfügung stehen.

\subsection{Items}
Items erfüllen die Rolle von Waffen und Belohnungen im Spiel, welche durch das Gewinnen von Minispielen oder Besiegen von Lehrern im Kampf gewonnen und im Kampf gegen selbige auch eingesetzt werden können. Jedes Item besitzt einen Namen, eine Bild, welches im Spiel angezeigt wird und einen Stärkegrad oder Level, welches den Schaden im Kampf determiniert. Es gibt 30 Items, sechs pro Fachschaft, mit jeweils zwei Items der Stärke Null, zwei der Stärke Eins und zwei der Stärke 2. Nach dem Abschluss eines Minispiels erhält der Spieler ein 
